\section{KL 梯度：相同积分下的不同控制}
\label{est:sec:kl-vector}
本节固定目标 $g=I_P(\ell\psi)$，总定理中的 $f_v=\ell_v\psi_v$ 是向量，随机估计量默认使用 $\beta=P$。数值估计量的分类逐项迁移到这个目标，但其控制函数和最优系数须在向量二阶矩下重新计算。沿用前节的 $h_g,\bar \psi_\tail,b_{\mathrm{TA}}$，并定义
\begin{equation*}
 \mu_g=\E_P[\ell_A\psi_A\mid J_A=1],\qquad
 d_{P_S}=\nabla P_S=\sum_{v\in S}P(v)\psi_v,\qquad g=h_g+\tau_P\mu_g.
\end{equation*}
头部截断梯度为 $h_g$，其偏差和方差为
\begin{equation}
 \operatorname{Bias}(h_g)=-\tau_P\mu_g,
 \qquad \Var(h_g)=0.
 \label{est:eq:chapter3-head-grad-bias}
\end{equation}

\subsection{HETS 与 ResOPD 的向量残差化}

尾部残差化在线策略蒸馏（On-Policy Distillation with Tail Residualization，ResOPD）从 HETS 的尾部梯度中减去可计算控制向量，再解析加回同一控制的期望。

HETS 梯度 estimator 定义为
\begin{equation}
 \widehat g^{\mathrm{HETS}}=h_g+J_A \ell_A\psi_A.
 \label{est:eq:chapter3-hets-grad}
\end{equation}
它无偏，且
\begin{equation}
 \Var(\widehat g^{\mathrm{HETS}})
 =\tau_P\Var_{\tail}(\ell_A\psi_A)+P_S\tau_P\lVert\mu_g\rVert_2^2.
 \label{est:eq:chapter3-hets-var}
\end{equation}

将尾部聚合为伪 token，定义
\begin{equation}
 b_{\mathrm{TA}}=\log\frac{\tau_P}{\tau_Q}\,\nabla_\theta\log\tau_P.
 \label{est:eq:chapter3-resopd-control}
\end{equation}
这个控制向量来自定理~\ref{inf:thm:kl-partition-gradient} 的粗粒度梯度。对头尾分区直接代入该定理，得到 $\nabla(H+\tau_P\ell_{\mathrm{mass}})=h_g+\tau_Pb_{\mathrm{TA}}$。因此完整梯度与粗粒度梯度之差为 $\tau_P(\mu_g-b_{\mathrm{TA}})$。使用 $J_A(\ell_A\psi_A-b_{\mathrm{TA}})$ 对该残差作无偏抽样，再加上粗粒度精确梯度，便得到 ResOPD。
对任意 action-independent 向量 $b$（与标量尾部形状项 $B$ 区分），定义统一控制族
\begin{equation}
 \widehat g(b)=h_g+\tau_Pb+J_A(\ell_A\psi_A-b).
 \label{est:eq:chapter3-control-family}
\end{equation}

\begin{theorem}[ResOPD 控制族的无偏性与方差]
\label{est:thm:chapter3-resopd}
在 $A\sim P$ 且 $b$ 在抽样前冻结时，
\begin{equation*}
 \E[\widehat g(b)]=g,
\end{equation*}
并且
\begin{equation}
 \Var(\widehat g(b))
 =\tau_P\Var_{\tail}(\ell_A\psi_A)
  +P_S\tau_P\lVert\mu_g-b\rVert_2^2.
 \label{est:eq:chapter3-control-var}
\end{equation}
HETS 是 $b=0$；ResOPD 是 $b=b_{\mathrm{TA}}$。
\end{theorem}

\begin{proof}
条件期望为
\begin{equation*}
 \E[\widehat g(b)]
 =h_g+\tau_Pb+\tau_P(\mu_g-b)=g.
\end{equation*}
给定 $J_A=0$ 时 estimator 为 $h_g+\tau_Pb$，给定 $J_A=1$ 时其条件均值为 $h_g+\tau_Pb+\mu_g-b$，两者均值差为 $\mu_g-b$。尾部条件下的随机部分为 $\ell_A\psi_A-b$，其协方差与 $\ell_A\psi_A$ 相同。对事件 $J_A$ 使用全协方差公式，得到式~\eqref{est:eq:chapter3-control-var}。
\end{proof}

ResOPD 与 HETS 的方差差为
\begin{equation}
 \Var(\widehat g^{\mathrm{HETS}})-\Var(\widehat g^{\mathrm{ResOPD}})
 =P_S\tau_P\left(2\langle\mu_g,b_{\mathrm{TA}}\rangle-\lVert b_{\mathrm{TA}}\rVert_2^2\right).
 \label{est:eq:chapter3-hets-resopd-gap}
\end{equation}
因此 ResOPD 的优势条件是
$\lVert\mu_g-b_{\mathrm{TA}}\rVert_2<\lVert\mu_g\rVert_2$。这正是 ResOPD 与 scalar KL 残差化的共同结构，区别只在于 scalar 的尾部均值是 $\mu_\ell$，向量梯度的尾部均值是 $\mu_g$。

该方差结论也有完整的协方差版本。记 $\Sigma_{\tail}=\operatorname{Cov}_{\tail}(\ell_A\psi_A)$，则
\begin{equation*}
 \operatorname{Cov}(\widehat g(b))
 =\tau_P\Sigma_{\tail}
  +P_S\tau_P(\mu_g-b)(\mu_g-b)^\top.
\end{equation*}
取迹即得到前面的方差公式。最优的无约束向量控制是 $b^*=\mu_g$，此时只剩 $\tau_P\Sigma_{\tail}$；该向量需要未知尾部，因此理论最优值不能作为 sparse payload 下的可计算算法。ResOPD 使用可计算的 $b_{\mathrm{TA}}$ 逼近它。

HETS 也不保证对每个分布都比直接 sampled-token 方差小。令 $\mu_{g,H}=\E[\ell_A\psi_A\mid A\in S]$、$V_H=\Var(\ell_A\psi_A\mid A\in S)$，全方差公式给出
\begin{align*}
 \Var(\ell_A\psi_A)&=P_SV_H+\tau_P\Var_{\tail}(\ell_A\psi_A)
 +P_S\tau_P\lVert\mu_{g,H}-\mu_g\rVert^2,\\
 \Var(\ell_A\psi_A)-\Var(\widehat g^{\mathrm{HETS}})
 &=P_SV_H+P_S\tau_P\{\lVert\mu_{g,H}\rVert^2
 -2\langle\mu_{g,H},\mu_g\rangle\}.
\end{align*}
头部精确求和删除了头部内部方差，但同时改变了事件间均值差，因此应使用这个差值判断排序。

\subsection{多样本、均方误差与支持集边界}

以上随机 estimator 均按单个动作计算。对固定 prefix 下 $m$ 个独立动作，使用同一确定性头部和控制项、对随机修正取平均，则偏差保持不变、协方差除以 $m$。标量均方误差为
\begin{equation*}
 \operatorname{MSE}(\widehat K)=\operatorname{Bias}(\widehat K)^2+\Var(\widehat K),
\end{equation*}
向量均方误差为
\begin{equation*}
 \E\lVert\widehat g-g\rVert^2
 =\lVert\operatorname{Bias}(\widehat g)\rVert^2+\Var(\widehat g).
\end{equation*}
因此确定性截断的零抽样方差不等于零估计误差；增加样本也无法消除它的偏差。这里的 $1/m$ 结论要求条件独立，同一条 trajectory 上不同 prefix 的 token 不能直接视为同一固定分布的独立重复。

当 $S=\mathcal V$ 时，按边界定义直接返回完整头部积分，不计算尾部 log-mass 或条件均值，尾部修正置零，方差为零；当 $S=\varnothing$ 时，HETS 退化为直接 sampled-token，$\tau_P=\tau_Q=1$ 使 ResOPD 控制为零。存在零概率时必须先检查绝对连续性及矩条件；尤其 K3 的 $\E_P[Q(A)/P(A)]=1$ 要求 $Q$ 的质量全部落在 $P$ 的支持上。
